\documentclass[10pt]{article}
\usepackage{template/template}
\usepackage{tikz}
\usetikzlibrary{positioning}
\NewsSetKeys{date={2026年11月21日},vol={039},title={每日有趣新闻},author={编辑部},note={资料窗口：11月19日至20日（北京时间）；本稿尚待构建、版式和原始图片验收，不是正式PDF。}}
\begin{document}
\begin{documents}
\newsitem{AI上市后的第二天：市场为何不肯只为增长故事买单？}{科技、金融与经济｜11月20日｜深度观察}{Anthropic于11月19日上市之后，全球科技股在20日继续承压。路透社报道，Anthropic股票盘中一度跌至接近发行价，纳斯达克指数当日大幅下挫；芯片、云计算及AI基础设施相关公司同步受到波及。\cite{reuters_ai,reuters_market} 这与昨日IPO首日的热烈交易形成鲜明对比，却不能单凭两日价格变化就宣布AI技术或商业模式失败。股价是市场对未来现金流、风险和折现率的综合报价，不是模型能力排行榜。

理解这次调整，首先要把模型公司的收入增长和资本市场对其估值分开。训练大模型需要昂贵的芯片、网络、电力、数据中心及研发人员，推理服务又有持续的边际成本。若收入增长的同时计算需求增长得更快，利润与现金流未必同步改善。与传统软件订阅不同，模型服务的单位计算成本可能随着模型复杂度、客户调用模式和竞争格局而变化，不能把软件业历史毛利率机械套用。

其次，AI产业链内部存在显著的时间错配。设备制造商可能在数据中心建成前确认订单，数据中心运营商可能在获得稳定客户前筹集资金，模型服务商则必须在收入真正实现前投入训练资源。对一家企业成立的增长叙事，并不自动意味着整个产业链的资本支出都能在预期年限内收回。投资者开始质疑的，是项目的成本、资金来源、最终付费客户和回收期限之间是否匹配。

利率预期也是不可忽略的一环。长期增长公司的估值对折现率尤其敏感：即使未来预期利润没有变化，市场要求的回报率上升，也会压低今天愿意支付的价格。与此同时，AI数据中心建设对电力和设备的集中需求可能在短期推高部分投入价格，而生产率收益未必马上兑现。芝加哥联储主席古尔斯比11月19日就提出AI投资可能造成短期供给约束和通胀压力的担忧；这是政策制定者的风险判断，不是已经证实的经济因果结论。\cite{reuters_fed}

因此，不能简单地把股票下跌解释成技术倒退，也不能以模型能力进步否定估值风险。更可靠的观察指标是未来季度的收入结构、客户留存、现金流、单位推理成本、数据中心利用率、融资条款及实际电力接入进度。对公众而言，AI产业的关键问题已经从“模型能做什么”逐步扩展为“谁为这些能力付费、何时付费，以及是否足以覆盖长期投入”。

本刊编辑判断：上市提供了更多可观察的财务披露机会，但首日上涨与次日回落都只是极短的市场样本。路透社的市场报道提供价格与交易背景，企业未来盈利能力仍需招股文件及后续财报验证。}
\newsitem{中美军用AI讨论：所谓“人类参与”究竟意味着什么？}{外交与技术治理｜11月20日｜重点报道}{美国财政部长贝森特11月20日表示，特朗普与习近平就重大军事行动中的人工智能使用讨论了人类参与问题，并称双方取得某种共识。\cite{reuters_military} 目前这一说法来自美方官员转述；在未见双方一致公布的协议文本之前，不能将其描述为已经生效、具备明确核查机制的双边军控条约。

“人在回路中”通常意味着机器提出建议或执行某些步骤，但人类保留关键决定权。然而，人类在什么阶段参与、是否能够真正理解系统输出、能否及时否决，以及自动化系统的速度是否使审查流于形式，都会改变这一原则的实际含义。对致命武器系统而言，目标识别、交战授权、战场通信和责任归属不能只用一句原则替代具体制度。

这次表态的重要性，在于军事AI的治理讨论可能进入最高层外交议程；它的局限，则是目前公开信息不足以确定适用武器类别、技术定义和执行机制。应继续核查中方正式通报、会谈纪要及后续国际谈判文件。}
\newsitem{先进芯片出口政策：一句“只卖旧芯片”距离法规有多远？}{中美经贸｜11月20日｜重点报道}{特朗普11月20日对记者表示，美国希望只向中国出售较旧的芯片，政府正在讨论一项协议，但未公布具体条款。\cite{reuters_chips} 这是总统公开表达的政策方向，不能直接等同于新的出口许可证已经签发、现行法规已经修改或企业已经获得确定的市场准入。

芯片管制不仅关乎产品型号，还涉及计算性能、互联能力、最终用户、转售地点及海外数据中心提供的计算服务。政策从政治声明转为实际规则，需要主管部门发布具体文本、定义适用范围及过渡安排。中国企业的替代研发与国际供应链调整也会影响政策效果，但不能从一次记者会直接推出结果。}
\newsitem{数学拾趣：一个反例如何改写图的分解阈值？}{数学｜11月20日arXiv新稿｜研究预印本}{Juho Lee在预印本《Star decomposition of complete multipartite graphs》中研究把图的边分解为大小为 $k$ 或 $k+1$ 的星形子图的问题。\cite{lee} 星形图 $K_{1,k}$ 由一个中心点与 $k$ 个叶子组成；所谓分解，是让每条边恰好属于一个指定类型的星形子图。

作者研究端点度数和 $\sigma(G)=\min_{uv\in E(G)}(d(u)+d(v))$ 对分解存在性的影响，并声称当 $k=3$ 时，完全多部图的最佳整数阈值为9；对一般 $k$，构造了三部图反例，说明所需阈值至少为 $k^2$。这里的困难不在于找一个局部星形，而在于所有边必须同时被恰好覆盖，局部可行性未必能拼成全局分解。作者使用有限计算验证与延拓论证处理特例，再用显式构造给出下界。

本文为2026年11月20日新提交的arXiv预印本。上述结论是作者在摘要中的主张，本刊未独立复核证明、计算程序及同行评议状态，不能称为已被学界确认的定理。}
\newsitem{航天与地球观测：卫星看见南太平洋风暴季}{NASA原始资料｜11月20日｜简讯}{NASA地球观测栏目11月20日发布南太平洋风暴影像与说明，介绍热带气旋Xavier开启当地风暴季的情况。\cite{nasa_storm} 卫星影像可以直接呈现云系结构，但云顶图像并不能单独提供地面风速、损失规模或长期气候趋势。正式报纸应采用NASA原始图像文件、注明观测日期、仪器与图像处理信息；本次尚未完成图片资源下载与归档，因此仅保留文字，不以AI合成图替代。}
\newsitem{图解：融资、估值与经营现金流}{编辑部自绘概念示意｜非按比例的数据图}{\begin{center}\begin{tikzpicture}[node distance=4mm, every node/.style={draw=newsblue,fill=newsblue!7,rounded corners,align=center,font=\scriptsize,minimum width=3.2cm,minimum height=7mm}]
\node (a) {融资所得现金\\一次性资本流入};
\node[below=of a] (b) {经营现金流\\持续收入减持续支出};
\node[below=of b] (c) {市场股权估值\\对未来收益的价格判断};
\end{tikzpicture}\end{center}
\footnotesize 编辑部原创分类示意，三项指标并非相同量，也不存在固定换算关系；不表示实际公司财务数据。}
\end{documents}
\small\section*{参考资料}
\bibliographystyle{plain}
\bibliography{sources/2026/Nov/21/refs}
\end{document}
